% Section structure adapted from 6.pdf; draft prompts are not thesis prose.
\chapter{Research Design and Measurement Methodology}
\label{ch:methodology}
\draftnote{Define comparison protocols once. Separate exploratory findings from questions specified before measurement; identify the actual evidence for each claim.}

\section{Research Questions as Controlled Comparisons}
\label{sec:3-1}
\draftnote{Write this section using the verified methods and evidence. State the relevant conditions and connect the argument to its research question.}

\section{Datasets, Queries, Ground Truth and Splits}
\label{sec:3-2}
\draftnote{Write this section using the verified methods and evidence. State the relevant conditions and connect the argument to its research question.}

\section{Measurement Boundary, Variance and Memory Accounting}
\label{sec:3-3}
\draftnote{Write this section using the verified methods and evidence. State the relevant conditions and connect the argument to its research question.}

\section{Baselines, Tuning Grids and Frontier Construction}
\label{sec:3-4}
\draftnote{Write this section using the verified methods and evidence. State the relevant conditions and connect the argument to its research question.}

\section{Versioned Claim-to-Evidence Ledger}
\label{sec:3-5}
\draftnote{Record commit, binary, compiler options, index/cache identity, query split, raw result, metric, and figure-generation script for every main claim.}

\section{Evidence Strength and Threat Controls}
\label{sec:3-6}
\draftnote{Write this section using the verified methods and evidence. State the relevant conditions and connect the argument to its research question.}
